\chapter{Experiment: Inspect CPU Execution and Stack Frames}

\section{Objective}

This experiment will help you understand how C programs execute at the hardware level by:
\begin{enumerate}
\item Examining generated assembly code to see CPU instructions
\item Observing stack frame layout and function call mechanics
\item Understanding register usage and calling conventions
\item Analyzing memory addresses and endianness
\end{enumerate}

\section{Setup Requirements}

\begin{itemize}
\item C compiler (gcc or clang)
\item Linux or WSL is recommended for consistent assembly output
\item Basic familiarity with command line tools
\item Debugger (gdb) optional but helpful
\end{itemize}

\section{Experiment 1: Generate and Inspect Assembly Code}

\subsection{Step 1: Create a Test Program}

Create a file called \texttt{add.c}:

\begin{lstlisting}[language=C]
int add(int a, int b) {
    int result = a + b;
    return result;
}

int main(void) {
    int sum = add(5, 3);
    return sum == 8 ? 0 : 1;
}
\end{lstlisting}

\subsection{Step 2: Generate Assembly Output}

Generate assembly code with frame pointers preserved (for clarity):

\begin{lstlisting}[language=bash]
gcc -S -fno-omit-frame-pointer -O0 -o add.s add.c
\end{lstlisting}

\subsection{Step 3: Examine the Assembly File}

Open \texttt{add.s} in your text editor:

\begin{lstlisting}[language=bash]
cat add.s
\end{lstlisting}

\textbf{What to look for:}

\begin{enumerate}
\item \textbf{Function prologue} in \texttt{add}:

\begin{lstlisting}
push   rbp            # Save caller's frame pointer
mov    rbp, rsp       # Establish new frame pointer
\end{lstlisting}

\item \textbf{Argument passing}: On x86-64 with System V ABI: \texttt{edi} contains first argument (a), \texttt{esi} contains second (b)

\item \textbf{Return value}: The sum is placed in \texttt{eax} before returning

\item \textbf{Function epilogue}:

\begin{lstlisting}
pop    rbp            # Restore caller's frame pointer
ret                   # Return to caller
\end{lstlisting}
\end{enumerate}

\subsection{Step 4: Compile and Run}

\begin{lstlisting}[language=bash]
gcc -o add add.c
./add
echo "Exit code: $?"
\end{lstlisting}

Expected exit code: 0 (success)

\section{Experiment 2: Examine Stack Frame Layout}

\subsection{Step 1: Create a Program with Local Variables}

Create \texttt{stack\_demo.c}:

\begin{lstlisting}[language=C]
int demonstrate_stack(int x) {
    int local1 = x * 2;
    int local2 = x + 10;
    return local1 + local2;
}

int main(void) {
    return demonstrate_stack(5);
}
\end{lstlisting}

\subsection{Step 2: Generate Assembly}

\begin{lstlisting}[language=bash]
gcc -S -fno-omit-frame-pointer -O0 -o stack_demo.s stack_demo.c
cat -n stack_demo.s | grep -A 20 "demonstrate_stack"
\end{lstlisting}

\textbf{What to observe:}

\begin{enumerate}
\item \textbf{Stack allocation}:

\begin{lstlisting}
sub    rsp, 16        # Allocate space for local variables
\end{lstlisting}

\item \textbf{Local variable access}: \texttt{local1} at \texttt{[rbp-4]}, \texttt{local2} at \texttt{[rbp-8]}
\end{enumerate}

\section{Experiment 3: Explore Endianness}

\subsection{Step 1: Create Endianness Test Program}

Create \texttt{endianness.c}:

\begin{lstlisting}[language=C]
#include <stdio.h>

int main(void) {
    unsigned int value = 0x12345678;
    unsigned char *bytes = (unsigned char *)&value;

    printf("Value: 0x%x\n", value);
    printf("Byte 0: 0x%02x\n", bytes[0]);
    printf("Byte 1: 0x%02x\n", bytes[1]);
    printf("Byte 2: 0x%02x\n", bytes[2]);
    printf("Byte 3: 0x%02x\n", bytes[3]);

    if (bytes[0] == 0x78) {
        printf("System is: Little-endian\n");
    } else if (bytes[0] == 0x12) {
        printf("System is: Big-endian\n");
    }

    return 0;
}
\end{lstlisting}

\subsection{Step 2: Compile and Run}

\begin{lstlisting}[language=bash]
gcc -o endianness endianness.c
./endianness
\end{lstlisting}

\textbf{Expected output on x86-64:}

\begin{lstlisting}
Value: 0x12345678
Byte 0: 0x78
Byte 1: 0x56
Byte 2: 0x34
Byte 3: 0x12
System is: Little-endian
\end{lstlisting}

\section{Experiment 4: Trace Register Usage with GDB}

\subsection{Step 1: Create a Debuggable Program}

Create \texttt{registers.c}:

\begin{lstlisting}[language=C]
int multiply(int a, int b) {
    return a * b;
}

int main(void) {
    int result = multiply(7, 6);
    return result == 42 ? 0 : 1;
}
\end{lstlisting}

\subsection{Step 2: Compile with Debug Symbols}

\begin{lstlisting}[language=bash]
gcc -g -o registers registers.c
\end{lstlisting}

\subsection{Step 3: Debug with GDB}

\begin{lstlisting}[language=bash]
gdb ./registers
\end{lstlisting}

In GDB, run these commands:

\begin{lstlisting}
(gdb) break main
(gdb) break multiply
(gdb) run
(gdb) info registers rdi rsi rdx rcx r8 r9
(gdb) step
(gdb) info registers rax rbx rcx rdx
(gdb) continue
(gdb) quit
\end{lstlisting}

\section{Experiment 5: Compare Optimized vs Unoptimized Code}

\subsection{Step 1: Create Optimization Test}

Create \texttt{optimization.c}:

\begin{lstlisting}[language=C]
int calculate(int a, int b, int c) {
    int temp1 = a + b;
    int temp2 = temp1 * c;
    int temp3 = temp2 - a;
    return temp3;
}

int main(void) {
    return calculate(5, 3, 2);
}
\end{lstlisting}

\subsection{Step 2: Generate and Compare}

\begin{lstlisting}[language=bash]
gcc -S -O0 -fno-omit-frame-pointer -o opt_unopt.s optimization.c
gcc -S -O2 -o opt_opt.s optimization.c
diff -y opt_unopt.s opt_opt.s | head -30
\end{lstlisting}

\section{Experiment 6: Examine Calling Convention with Many Arguments}

\subsection{Step 1: Create Multi-Argument Function}

Create \texttt{many\_args.c}:

\begin{lstlisting}[language=C]
long many_arguments(long a, long b, long c, long d,
                    long e, long f, long g, long h) {
    return a + b + c + d + e + f + g + h;
}

int main(void) {
    return (int)many_arguments(1, 2, 3, 4, 5, 6, 7, 8);
}
\end{lstlisting}

\subsection{Step 2: Generate Assembly}

\begin{lstlisting}[language=bash]
gcc -S -O0 -fno-omit-frame-pointer -o many_args.s many_args.c
cat -n many_args.s | grep -A 30 "many_arguments:"
\end{lstlisting}

\textbf{What to observe (System V AMD64 ABI):}

Arguments 1-6: Passed in registers (rdi, rsi, rdx, rcx, r8, r9)

Arguments 7-8: Passed on the stack

\section{Questions for Reader}

\begin{enumerate}
\item Which register holds the return value for integer functions on x86-64?
\item Does the stack grow toward higher or lower addresses?
\item What is the purpose of \texttt{rbp}?
\item In the System V AMD64 ABI, where are the first 6 integer arguments passed?
\item How does optimization level (-O0 vs -O2) affect the generated assembly?
\end{enumerate}

\section{Summary}

Through these experiments, you should now understand:

\begin{itemize}
\item \textbf{CPU Execution Model}: How the fetch-decode-execute cycle works
\item \textbf{Register Usage}: How registers are used for arguments, return values, and temporary storage
\item \textbf{Stack Mechanics}: How stack frames are created and destroyed
\item \textbf{Calling Conventions}: The rules that govern function calls
\item \textbf{Memory Organization}: How data is stored in memory (endianness)
\end{itemize}
